% GENERATED FILE. Edit canonical lessons, then run npm run generate:books.
% canonical-source-hash: fnv1a64:d67e297ad1c2b19d
% canonical-lessons: GE-C409-baustelle, GE-C409-umleitung, GE-C409-umleiten, GE-C409-sperren, GE-C409-raststaette

\chapter{Bauarbeiten}
\label{ch:ge-bauarbeiten}
\hlchaptermodality{409}

\begin{chapteropening}
\textbf{By the end of this chapter:} \emph{I can follow road works and diversions: the road works, the diversion, diverting traffic, closing a road and the service area.}

Complete the last lesson of chapter 409: I can follow road works and diversions: the road works, the diversion, diverting traffic, closing a road and the service area.
\end{chapteropening}

\section[die Baustelle]{die Baustelle — the road works; the building site}

\label{lesson:GE-C409-baustelle}



\begin{quote}
Before the new one: say the German for the oncoming traffic, then the German for to swerve.
\end{quote}

\subsection*{You'll want to know: die Baustelle}
\textbf{die Baustelle} — \textquotedblleft{}the road works; the building site\textquotedblright{}.

\begin{quote}
Wegen der Baustelle fährt der Bus heute langsam.
\end{quote}

\begin{quote}
\emph{Because of the road works the bus is slow today.}
\end{quote}

Plural: \emph{die Baustellen}.

\begin{cousinweb}
\emph{Baustelle} is built from \emph{Bau} 'building', from \emph{bauen} 'to build', inherited from Proto-Germanic \emph{būaną} (EtymDB), and \emph{Stelle} 'place'.
\end{cousinweb}

\begin{grammarlens}[title={What you've built}]
One more word for telling, explaining and describing at B1.
\end{grammarlens}

\subsection*{Your turn}
\begin{itemize}
\raggedright
  \item \emph{Say it:} \emph{die Baustelle}
  \item \emph{Say it:} \emph{die Baustelle}, once more
  \item \emph{Say it:} a sentence of your own with \emph{die Baustelle}
  \item \emph{Recall:} say the German for the oncoming traffic, then the German for to swerve, then say \emph{die Baustelle} again
\end{itemize}

\begin{tcolorbox}[breakable,colback=teal!4,colframe=teal!35!black,title={Before you move on}]
What does die Baustelle mean? (\textquotedblleft{}The road works; the building site\textquotedblright{}.) Say it once more.
\end{tcolorbox}
\section[die Umleitung]{die Umleitung — the diversion (US: detour)}

\label{lesson:GE-C409-umleitung}



\begin{quote}
Before the new one: say the German for to swerve, then the German for the road works.
\end{quote}

\subsection*{You'll want to know: die Umleitung}
\textbf{die Umleitung} — \textquotedblleft{}the diversion (US: detour)\textquotedblright{}.

\begin{quote}
Die Umleitung führt durch drei kleine Dörfer.
\end{quote}

\begin{quote}
\emph{The diversion goes through three small villages.}
\end{quote}

Plural: \emph{die Umleitungen}.

\begin{cousinweb}
\emph{Umleitung} is built from \emph{um} 'round, around' and \emph{leiten} 'to lead, to guide', inherited from Old High German \emph{leiten} and a cousin of English \emph{lead} (EtymDB).
\end{cousinweb}

\begin{grammarlens}[title={What you've built}]
One more word for telling, explaining and describing at B1.
\end{grammarlens}

\subsection*{Your turn}
\begin{itemize}
\raggedright
  \item \emph{Say it:} \emph{die Umleitung}
  \item \emph{Say it:} \emph{die Umleitung}, once more
  \item \emph{Say it:} a sentence of your own with \emph{die Umleitung}
  \item \emph{Recall:} say the German for to swerve, then the German for the road works, then say \emph{die Umleitung} again
\end{itemize}

\begin{tcolorbox}[breakable,colback=teal!4,colframe=teal!35!black,title={Before you move on}]
What does die Umleitung mean? (\textquotedblleft{}The diversion (US: detour)\textquotedblright{}.) Say it once more.
\end{tcolorbox}
\section[umleiten]{umleiten — to divert (traffic)}

\label{lesson:GE-C409-umleiten}



\begin{quote}
Before the new one: say the German for the road works, then the German for the diversion.
\end{quote}

\subsection*{You'll want to know: umleiten}
\textbf{umleiten} — \textquotedblleft{}to divert (traffic)\textquotedblright{}.

\begin{quote}
Die Polizei leitet die Autos über die Landstraße um.
\end{quote}

\begin{quote}
\emph{The police divert the cars along the country road.}
\end{quote}

Separable: \emph{sie leitet um}, \emph{sie hat umgeleitet}.

\begin{cousinweb}
\emph{umleiten} is built from \emph{um} 'round' and \emph{leiten} 'to lead', an inherited verb from Old High German \emph{leiten} (EtymDB).
\end{cousinweb}

\begin{grammarlens}[title={What you've built}]
One more word for telling, explaining and describing at B1.
\end{grammarlens}

\subsection*{Your turn}
\begin{itemize}
\raggedright
  \item \emph{Say it:} \emph{umleiten}
  \item \emph{Say it:} \emph{umleiten}, once more
  \item \emph{Say it:} a sentence of your own with \emph{umleiten}
  \item \emph{Recall:} say the German for the road works, then the German for the diversion, then say \emph{umleiten} again
\end{itemize}

\begin{tcolorbox}[breakable,colback=teal!4,colframe=teal!35!black,title={Before you move on}]
What does umleiten mean? (\textquotedblleft{}To divert (traffic)\textquotedblright{}.) Say it once more.
\end{tcolorbox}
\section[sperren]{sperren — to close (a road); to block}

\label{lesson:GE-C409-sperren}



\begin{quote}
Before the new one: say the German for the diversion, then the German for to divert.
\end{quote}

\subsection*{You'll want to know: sperren}
\textbf{sperren} — \textquotedblleft{}to close (a road); to block\textquotedblright{}.

\begin{quote}
Nach dem Unfall sperrt die Polizei die Brücke.
\end{quote}

\begin{quote}
\emph{After the accident the police close the bridge.}
\end{quote}

\emph{Gesperrt} — closed, on signs. \emph{Eine Karte sperren} — to block a card.

\begin{cousinweb}
\emph{sperren} is an inherited Germanic verb; EtymDB links it with English \emph{spar} 'to bolt, to bar', so it is usually explained as first meaning to shut something with a bar.
\end{cousinweb}

\begin{grammarlens}[title={What you've built}]
One more word for telling, explaining and describing at B1.
\end{grammarlens}

\subsection*{Your turn}
\begin{itemize}
\raggedright
  \item \emph{Say it:} \emph{sperren}
  \item \emph{Say it:} \emph{sperren}, once more
  \item \emph{Say it:} a sentence of your own with \emph{sperren}
  \item \emph{Recall:} say the German for the diversion, then the German for to divert, then say \emph{sperren} again
\end{itemize}

\begin{tcolorbox}[breakable,colback=teal!4,colframe=teal!35!black,title={Before you move on}]
What does sperren mean? (\textquotedblleft{}To close (a road); to block\textquotedblright{}.) Say it once more.
\end{tcolorbox}
\section[die Raststätte]{die Raststätte — the motorway service area}

\label{lesson:GE-C409-raststaette}



\begin{quote}
Before the new one: say the German for to divert, then the German for to close.
\end{quote}

\subsection*{You'll want to know: die Raststätte}
\textbf{die Raststätte} — \textquotedblleft{}the motorway service area\textquotedblright{}.

\begin{quote}
An der Raststätte trinken wir einen Kaffee.
\end{quote}

\begin{quote}
\emph{At the service area we have a coffee.}
\end{quote}

Plural: \emph{die Raststätten}. A smaller one without a restaurant is \emph{der Rastplatz}.

\begin{cousinweb}
\emph{Raststätte} is built from \emph{Rast} 'rest, a break on a journey' and \emph{Stätte} 'place', inherited from Middle High German \emph{stat} (EtymDB).
\end{cousinweb}

\begin{grammarlens}[title={What you've built}]
One more word for telling, explaining and describing at B1.
\end{grammarlens}

\subsection*{Your turn}
\begin{itemize}
\raggedright
  \item \emph{Say it:} \emph{die Raststätte}
  \item \emph{Say it:} \emph{die Raststätte}, once more
  \item \emph{Say it:} a sentence of your own with \emph{die Raststätte}
  \item \emph{Recall:} say the German for to divert, then the German for to close, then say \emph{die Raststätte} again
\end{itemize}

\begin{tcolorbox}[breakable,colback=teal!4,colframe=teal!35!black,title={Before you move on}]
What does die Raststätte mean? (\textquotedblleft{}The motorway service area\textquotedblright{}.) Say it once more.
\end{tcolorbox}
